\documentclass[10pt]{article}
\usepackage[margin=0.9in]{geometry}\usepackage{amsmath,amssymb,booktabs,microtype,xcolor,fancyhdr}\usepackage[hidelinks]{hyperref}
\definecolor{navy}{RGB}{24,55,95}\pagestyle{fancy}\fancyhf{}\lhead{$E_7$ on $V_{56}$}\rhead{Proof and verification strategy}\cfoot{\thepage}\setlength{\headheight}{14pt}
\newcommand{\Sg}{\mathcal S}\newcommand{\CT}{\operatorname{CT}}
\title{\color{navy}\textbf{The compact $E_7$ sigma-MGF on $V_{56}$}\\[1mm]\large Proof and dimension-aware verification strategy}
\author{Computational proof companion}\date{17 July 2026}
\begin{document}\maketitle
\begin{abstract}
For simply connected compact $E_7$ on its minuscule 56, we prove the
Molien--Weyl expression and parity rule, exactly verify the full quadratic
expansion with explicit torus matrices, and state the degree-four expansion
using cited invariant theory.  We also give a sparse quartic verification
algorithm; the present artifact does not report that computation as executed.
\end{abstract}

\section{Formula and parity}
The Cauchy identity and Haar projection prove
\[
 \Sg_{K,V}=\int_K\prod_a\det(I-t_a\rho(g))^{-1}dg
 =\sum_\lambda\dim(\mathbb S_\lambda V)^K s_\lambda(\mathbf t).
\]
The equivalent monomial coefficient is
$[m_\tau]\Sg_{K,V}=\dim(\bigotimes_i\operatorname{Sym}^{\tau_i}V)^K$.
The 56 weights of $V_{56}$ form one Weyl orbit.  Hence
\[
 \boxed{\Sg_{E_7,56}=\frac1{2903040}\CT_z\left[D_{E_7}(z)
 \prod_a\prod_{\mu\in W\omega_7}(1-t_a z^\mu)^{-1}\right].}
\]
The nontrivial central element acts as $-I$, so every odd-total-degree
coefficient vanishes.  This statement concerns simply connected $E_7$; the
56 does not descend to $E_7/Z_2$.

\section{Low-degree tensors}
Let $\omega\in\bigwedge^2V^*$ be the invariant symplectic form and let
$q\in\operatorname{Sym}^4V^*$ be the Cartan--Freudenthal quartic.  At degree
four, the three pairwise contractions made from two copies of $\omega$
decompose under $S_4$ as the Schur types $(2,2)$ and $(1,1,1,1)$.  The quartic
$q$ supplies type $(4)$.  The one-copy statement
$\mathbb C[V_{56}]^{E_7}=\mathbb C[q]$ and the exhaustion of degree-four
tensor invariants by the symplectic pairings and $q$ are external
invariant-theory inputs; see Brown \cite{brown}.  Using those inputs,
\[
 \boxed{\Sg_{E_7,56}=1+e_2+h_4+s_{(2,2)}+e_4+O_{\ge6}},
 \qquad \boxed{\Sg(t,0,\ldots)=\frac1{1-t^4}}.
\]
The one-variable specialization kills $e_2,s_{(2,2)},e_4$ and retains the
quartic generator.
In particular $([m_{(2)}],[m_{(1,1)}])=(0,1)$, while the predicted degree-four
row is
\[
([m_{(4)}],[m_{(3,1)}],[m_{(2,2)}],[m_{(2,1,1)}],[m_{(1,1,1,1)}])
=(1,1,2,2,4).
\]

\section{Executed quadratic verification and optional quartic confirmation}
The notebook generates the 126 roots and the 56-weight orbit, constructs
$M(\theta)=\operatorname{diag}(e^{i\langle\mu,\theta\rangle})$, and verifies
the determinant integrand.  Jacobi--Trudi characters are integrated exactly
using
\[
 \frac1{|W|}\CT(Df)=\CT\left(f\prod_{\alpha>0}(1-z^{-\alpha})\right).
\]
The complete quadratic computation is
\begin{center}\small
\begin{tabular}{@{}c|rrrr@{}}\toprule
$\lambda$&$\varnothing$&$(1)$&$(2)$&$(1,1)$\\\midrule
proposed&1&0&0&1\\
exact Weyl check&1&0&0&1\\\bottomrule
\end{tabular}
\end{center}
All entries pass.  The matrix unitarity residual is $4.2\times10^{-16}$ and
the determinant/eigenvalue residual is $8.6\times10^{-14}$.

\paragraph{Status of the degree-four row.}
The row follows from the cited invariant-theory decomposition.  The notebook
output displayed here verifies only degrees at most two.  A fully mechanical
confirmation may be obtained as follows, but is not executed in this
artifact.  Do not form dense matrices on
$V^{\otimes4}$ (dimension $56^4$).  Instead: (i) generate the weight-sparse
root-operator action on each of $\operatorname{Sym}^4V$, $\mathbb S_{(2,2)}V$,
and $\bigwedge^4V$; (ii) retain only zero-weight basis vectors; (iii) solve the
sparse raising-operator nullspace; and (iv) check dimensions $1,1,1$.
Equivalently, stream the three Schur characters through the same exact Weyl
functional.  An executed result must report the zero-weight dimensions,
exact raising-kernel ranks, arithmetic mode, and output hash.

\paragraph{Reproducibility.} Run \texttt{marimo\_notebooks/lie\_group\_e7.py} with marimo.
\begin{thebibliography}{9}
\bibitem{brown} R.~B.~Brown, ``Groups of type $E_7$,'' \emph{J. Reine Angew.
Math.} \textbf{236} (1969), 79--102.
\bibitem{adams} J.~F.~Adams, \emph{Lectures on Exceptional Lie Groups},
Chicago Lectures in Mathematics (1996).
\end{thebibliography}
\end{document}
